% !TEX root = ../main.tex
\chapter{Apple Silicon Architecture}
\label{ch:apple_silicon}

The substrate on which \Lucid\ executes is not a generic compute fabric.
It is a tightly integrated, vertically designed system-on-chip family
(\emph{Apple silicon}, sometimes \emph{Apple~M}) that Apple introduced in
November 2020 with the M1 and has iterated on each subsequent year through
the M2 (2022), M3 (2023), and M4 (2024) generations. Every architectural
choice in \Lucid\ --- from the device split into a CPU stream and a GPU
stream (\chref{ch:backend_dispatch}), to the carve-out that runs linear
algebra through \MLX\ even when the user selected the CPU
(\chref{ch:linalg}), to the deliberate avoidance of the Neural Engine
(\secref{sec:apple_silicon:ane}) --- presupposes a reader who understands
the chip beneath the runtime. This chapter exists to provide that
understanding.

We approach the silicon the way a compiler writer would: as a
heterogeneous collection of compute units sharing a single unified memory
through a coherent fabric, with each unit exposing a distinct programming
model and a distinct cost profile.\footnote{We deliberately defer
\emph{unified memory} itself to \chref{ch:unified_memory}; here we treat
it as a property of the SoC that gives every compute unit a coherent view
of the same DRAM pool, and we return to its detailed model later.} The
exposition is organised by compute unit. \secref{sec:apple_silicon:cpu}
covers the CPU complex, with its heterogeneous performance and efficiency
cores. \secref{sec:apple_silicon:amx} treats the Apple Matrix Coprocessor
(AMX) and its transition to the ARM Scalable Matrix Extension (SME) in
the M4. \secref{sec:apple_silicon:ane} discusses the Neural Engine and
explains, with reference to \Lucid's design constraints, why it is
intentionally absent from \Lucid's dispatch graph.
\secref{sec:apple_silicon:gpu} describes the GPU, which is the workhorse
of nearly every training workload on the platform.
\secref{sec:apple_silicon:memory} addresses the memory hierarchy and
system-level cache (SLC). Finally,
\secref{sec:apple_silicon:generations} compares the four generations
quantitatively, and \secref{sec:apple_silicon:ml_implications} draws the
conclusions that motivate the rest of the book.

\section{Overview}
\label{sec:apple_silicon:overview}

Apple silicon refers to a family of ARMv8 (and, from the M4 onwards,
ARMv9-A) SoCs designed by Apple's silicon engineering group and
manufactured by TSMC on its leading-edge processes
\cite{apple_m1_2020,apple_m4_2024}. The first member,
M1, shipped in late 2020 on TSMC's first-generation 5\,nm node (N5);
M2 used an enhanced 5\,nm node (N5P) in mid-2022; M3 moved to a
first-generation 3\,nm node (N3B) in late 2023; M4 moved again to
TSMC's N3E in mid-2024.\footnote{Process generation matters more than
it might appear: each shrink alters not only the area budget but the
voltage/frequency curve, which feeds into the power envelope that
ultimately bounds sustained ML throughput. We revisit this in
\secref{sec:apple_silicon:ml_implications}.}

The hallmark of Apple silicon, and the property that most distinguishes
it from contemporaneous discrete-GPU systems, is the integration of
heterogeneous compute units --- a CPU complex, a GPU, a Neural Engine
(NPU), and one or more matrix coprocessors --- on a single die, all
accessing a single physical DRAM pool through a single coherent fabric.
There is no PCIe boundary between the CPU and the accelerator; there is
no explicit \code{cudaMemcpy} step; the GPU does not have a private
VRAM. This is what \emph{unified memory} means at the SoC level, and it
is the architectural reason why \Lucid's CPU--GPU bridge is designed
around zero-copy view sharing rather than data marshalling
(\chref{ch:unified_memory}, \chref{ch:bridge_boundaries}).

The family has four published tiers, each obtained by stamping
additional copies of the same building blocks:
\begin{itemize}
  \item \textbf{Base} (M1 / M2 / M3 / M4): one CPU complex with a small
    P-core cluster (typically 4 cores) and a small E-core cluster
    (typically 4 cores), a moderate GPU (7--10 cores), 8--24\,GB DRAM,
    and roughly 60--120\,GB/s of memory bandwidth.
  \item \textbf{Pro}: an enlarged P-core cluster (6--10 cores), a
    larger GPU (14--20 cores), wider memory bus (4 channels of
    LPDDR5), and 200\,GB/s-class bandwidth.
  \item \textbf{Max}: another doubling of the P-core complex and GPU
    (24--40 GPU cores), 8 channels of LPDDR5, and 400\,GB/s-class
    bandwidth.
  \item \textbf{Ultra}: two Max dies connected through Apple's
    \emph{UltraFusion} silicon interposer to present a single logical
    SoC with up to 64 GPU cores and 800\,GB/s of aggregate bandwidth.
\end{itemize}

For \Lucid, the relevant generalisation is the following: as one moves
from base to Ultra, the available GPU compute grows roughly $8\times$,
memory bandwidth grows roughly $13\times$, but per-thread CPU
performance changes only marginally. This shapes the dispatch
philosophy laid out in \chref{ch:backend_dispatch}: any
performance-critical work that can be expressed as a GPU primitive
must be sent to \MLX, because the CPU side does not scale with the
chip tier.

\begin{figure}[t]
\centering
\begin{adjustbox}{max width=\textwidth, center}
\begin{tikzpicture}[
  block/.style={draw, rounded corners=2pt, minimum height=22pt,
                minimum width=2.6cm, font=\small\sffamily,
                inner sep=4pt, align=center},
  cpu/.style={block, fill=lucid.note.bg},
  gpu/.style={block, fill=lucid.info.bg},
  acc/.style={block, fill=lucid.code.bg},
  fabric/.style={draw, fill=lucid.rule, minimum width=14cm,
                 minimum height=18pt, font=\small\sffamily\bfseries,
                 align=center},
  mem/.style={draw, rounded corners=2pt, fill=lucid.ok.bg,
              minimum height=22pt, minimum width=2.8cm,
              font=\small\sffamily, align=center, inner sep=4pt},
  dram/.style={draw, dashed, minimum width=14cm, minimum height=20pt,
               font=\small\sffamily\itshape, align=center},
  link/.style={thick, color=lucid.muted},
]
  % Compute tier (top)
  \node[cpu] (pc)    at (-5.5, 2.4) {P-cores};
  \node[acc] (amx)   at (-2.8, 2.4) {AMX / SME};
  \node[cpu] (ec)    at ( 0,   2.4) {E-cores};
  \node[gpu] (gpu)   at ( 2.8, 2.4) {GPU cores};
  \node[gpu] (ane)   at ( 5.5, 2.4) {Neural Engine};

  % Fabric
  \node[fabric] (fab) at (0, 0.3) {Apple Fabric (coherent NoC)};

  % Memory tier (bottom)
  \node[mem] (slc) at (-4, -1.6) {System Level\\Cache};
  \node[mem] (mc)  at ( 4, -1.6) {Memory\\Controllers};
  \node[dram] (dram) at (0, -3.3) {LPDDR4X / LPDDR5(X) DRAM};

  % Vertical links from compute to fabric
  \draw[link] (pc.south)  -- (pc.south  |- fab.north);
  \draw[link] (amx.south) -- (amx.south |- fab.north);
  \draw[link] (ec.south)  -- (ec.south  |- fab.north);
  \draw[link] (gpu.south) -- (gpu.south |- fab.north);
  \draw[link] (ane.south) -- (ane.south |- fab.north);

  % Vertical links from fabric to memory
  \draw[link] (slc.north) -- (slc.north |- fab.south);
  \draw[link] (mc.north)  -- (mc.north  |- fab.south);

  % SLC + MC to DRAM
  \draw[link] (slc.south) -- (slc.south |- dram.north);
  \draw[link] (mc.south)  -- (mc.south  |- dram.north);
\end{tikzpicture}
\end{adjustbox}
\caption[Apple silicon SoC topology.]{Schematic floorplan of an Apple
silicon SoC. All compute units --- the CPU complex (P-cores, E-cores,
AMX/SME), the GPU, the Neural Engine, and fixed-function media blocks
--- communicate with DRAM through a single coherent fabric (the Apple
Fabric, sometimes called the system-level interconnect) that also
hosts the System Level Cache. \Lucid\ targets the CPU complex (through
\Accelerate) and the GPU (through \MLX\ and \MPSGraph); the Neural
Engine is intentionally not part of the dispatch graph, for reasons
described in \secref{sec:apple_silicon:ane}.}
\label{fig:apple_silicon_topology}
\end{figure}

The remainder of this chapter inspects each block in
\figref{fig:apple_silicon_topology} in turn. We pay particular attention
to (i)~the contracts the unit exposes to a runtime, (ii)~the
microarchitectural quirks that shape \Lucid's kernels and dispatch
policy, and (iii)~the generation-over-generation drift that affects
backward compatibility and version-gated code paths.

\section{SoC Floorplan and Coherent Fabric}
\label{sec:apple_silicon:floorplan}

A modern Apple SoC is not a monolithic die in the architectural sense
that an Intel Xeon of comparable transistor count is. It is, instead,
a tiled die: a grid of physically distinct \emph{clusters} (CPU
clusters, GPU cluster groups, the Neural Engine, the media engines,
and so on) connected by a network-on-chip (NoC) that Apple internally
calls the \emph{Apple Fabric}. The fabric is a coherent interconnect:
loads and stores from any client (a CPU core, a GPU shader core, the
Neural Engine) observe a single, well-ordered view of memory, with the
fabric and the System Level Cache acting in concert to ensure that a
write by the GPU becomes visible to a subsequent CPU load without an
explicit invalidation.

This coherence guarantee is what enables \Lucid's bridge boundary
(\chref{ch:bridge_boundaries}) to be a metadata operation rather than
a \code{memcpy}: when \Lucid\ converts an \MLX\ array to a
contiguous CPU view, the underlying bytes are not moved, only
relabelled.\footnote{The mechanism is described in detail in
\chref{ch:tensor_model} and the \texttt{SharedStorage} class in
\code{lucid/\_C/backend/gpu/SharedStorage.cpp}. The short version is
that we extract the \code{MTL::Buffer*} backing an \code{mlx::core::array}
through \code{mlx::core::array::data<T>()}, then expose its bytes
through the same Accelerate-facing pointer that the CPU backend would
have used.}

The fabric has three properties that matter for ML workloads.

\paragraph{Coherence at the cache-line granularity.} Coherence is
managed at the standard ARM cache-line width (128 bytes on Apple
silicon's L1 and L2; the SLC operates at the same granularity to the
fabric). This means that false-sharing pathologies that plague
multi-socket x86 systems --- where two cores writing into adjacent
cache lines stall each other --- still exist on Apple silicon, but
they are bounded to the local cluster, since the fabric is the only
coherence point. In practice this affects \Lucid\ only in the
parameter-update path of distributed training (which is not yet
supported; see \chref{ch:reproducibility}).

\paragraph{Cluster-local L2 caches.} Each CPU cluster owns a private
L2 (typically 12--16\,MB for the P-cluster and 4\,MB for the
E-cluster). Inter-cluster traffic --- including all traffic from CPU
to GPU --- must traverse the fabric and is filtered by the SLC.
Latency to the SLC is roughly $80$ to $120$ cycles depending on
generation; latency to DRAM, when SLC misses, is in the $150$--$200$
cycle range on the P-cores. These numbers are not published by Apple
and have been characterised by independent researchers
\cite{johnson_m1_microarch_2021}.\footnote{The
most thorough published characterisation of the M1 microarchitecture
is by Dougall Johnson and the Asahi Linux project. Subsequent
generations have been benchmarked by SemiAnalysis and Anandtech-style
deep dives; numbers in this chapter are typical of those reports and
should be treated as approximate.}

\paragraph{Asymmetric bandwidth.} The fabric is not bandwidth-uniform.
The GPU has a wider connection to the fabric than any single CPU
cluster; the Neural Engine has its own dedicated path. In aggregate,
the GPU can saturate roughly $80$--$90\,\%$ of the chip's memory
bandwidth on a sustained streaming workload, while a single CPU core
can drive perhaps $30$--$40\,\%$.\footnote{This asymmetry is part of
why \Lucid's CPU backend uses \BNNS\ for convolutions: it is the only
CPU-side API on macOS that accesses the fabric through a
matrix-coprocessor-accelerated path rather than a single-core path.}
The Ultra variants double the fabric and thus the aggregate
bandwidth, but the per-client peak does not scale linearly.

\section{The CPU Complex}
\label{sec:apple_silicon:cpu}

Apple's CPU complex is a heterogeneous (\emph{big.LITTLE}) design that
combines a small number of large, out-of-order \emph{performance
cores} (P-cores) with a larger number of small, energy-optimised
\emph{efficiency cores} (E-cores). The cluster organisation has
remained stable across generations, but every core microarchitecture
has been re-implemented at least once.

\subsection{Performance cores}
\label{ssec:apple_silicon:p_cores}

The P-core is the most-quoted aspect of Apple silicon's
microarchitecture, and the most thoroughly reverse-engineered. Its
distinguishing properties are well-documented through Apple's
patents, Asahi Linux's reverse engineering, and Dougall Johnson's
public micro-benchmarks. We summarise the salient design choices.

\subsubsection{Pipeline width and out-of-order resources}

Every Apple P-core since Firestorm (M1) decodes \emph{eight
instructions per cycle} and dispatches across roughly the same number
of issue ports. This is the widest sustained-decode design in any
shipping general-purpose CPU and roughly twice the width of
contemporaneous x86 cores (which decode 4--6 instructions per cycle).
The implication is that ILP-rich code --- pointer chasing, tight
arithmetic loops --- runs unusually fast per cycle on Apple P-cores,
which is why ARM-translated Intel binaries (\code{Rosetta 2}) reach
$60$--$80\,\%$ of native performance despite the translation overhead.

The reorder buffer (ROB) on Firestorm is roughly $630$ entries; later
generations have grown this further. For comparison, contemporaneous
x86 ROBs are in the $350$--$500$ range. The wide front-end and deep
ROB are what enable the P-core to absorb the latency to the SLC and
DRAM without stalling.

\subsubsection{NEON and the SIMD path}

Each P-core has four 128-bit NEON pipes; this corresponds to four
single-precision multiply--add throughputs per cycle per core, or
$128$ GFLOP/s at 4\,GHz. NEON is the SIMD ISA \Accelerate's
\code{vDSP} layer targets directly; the \code{vDSP\_vmul}, \code{vsmsa},
and similar primitives in \Lucid's CPU backend
(\chref{ch:backend_dispatch}) compile down to NEON.

Crucially, NEON is \emph{not} the path used by \BNNS\ for
convolutions. \BNNS\ goes through the AMX coprocessor
(\secref{sec:apple_silicon:amx}), which is shared by the whole
P-cluster and can sustain matrix-multiply throughput an order of
magnitude higher than four NEON pipes can.

\subsubsection{SME and the M4 transition}

The M4 generation is the first Apple silicon to expose the ARM
Scalable Matrix Extension (SME) \cite{arm_sme_2021} at the ISA
level. SME is the architected successor to the (still
undocumented) AMX instruction set;
it gives software a portable, ARM-blessed API for matrix-multiply
acceleration. SME operates on \emph{tiles} of variable width
(\code{ZA} register file) and provides outer-product accumulation
instructions analogous to AMX's \code{outer-product} micro-ops.

For \Lucid\ the practical implication is small in the short term and
large in the long term:
\begin{enumerate}
  \item In the short term, \Accelerate's \BNNS\ on M4 transparently
    uses SME where it would have used AMX on M1--M3, so \Lucid's CPU
    backend gets the benefit without any code change.
  \item In the long term, SME's portability means that \Lucid\ could
    eventually emit SME intrinsics directly from \code{lucid.compile}
    (\chref{ch:compile_design}) without depending on \Accelerate. This
    is on the roadmap but is not in the 3.5 series.
\end{enumerate}

\subsubsection{Performance counters and observation}

The P-cores expose roughly $400$ performance monitoring unit (PMU)
events. These are not exposed through standard ARM PMU mechanisms on
macOS; instead, the \code{Instruments.app} CPU Counters template (or
the underlying \code{kperf} API) is the only supported access path.
\Lucid's profiler (\chref{ch:profiler_determinism}) uses
\code{mach\_absolute\_time()} for wall-clock timing and does not
currently read PMU counters; the user is expected to use Instruments
for cycle-level analysis.

\subsection{Efficiency cores}
\label{ssec:apple_silicon:e_cores}

The E-cores are interesting for \Lucid\ less because of what they
execute than because of what they \emph{do not} execute. They are
designed for sustained low-power operation, with a much narrower
front-end (typically 4-wide decode) and a smaller ROB ($\sim 180$
entries on Icestorm). They run at lower frequencies (peak roughly
$2$\,GHz, against $3.2$--$4.5$\,GHz for the P-cores) and lack the
parallel NEON throughput of the P-cores.

For an ML workload, the consequence is that any thread which falls
onto an E-core will be roughly $4$--$6\times$ slower per cycle than
the same thread on a P-core, and run at a lower clock besides. This
matters because the macOS scheduler will move threads off the P-cores
when their QoS class is set to background or utility
(\secref{ssec:apple_silicon:scheduling}). \Lucid's main training
loop, by design, uses no thread-affinity hints and runs at the
default QoS class \code{USER\_INITIATED}; this places it on P-cores
unless the user has dropped to background, which is the correct
behaviour for training workloads but means that \Lucid\ will
\emph{not} oversubscribe E-cores even when they are idle.

\subsection{Big.LITTLE scheduling and QoS}
\label{ssec:apple_silicon:scheduling}

macOS implements scheduling through Grand Central Dispatch (GCD),
which expresses thread priority through QoS classes:
\code{USER\_INTERACTIVE}, \code{USER\_INITIATED}, \code{UTILITY},
\code{BACKGROUND}, and \code{DEFAULT}. The lower three classes
(\code{UTILITY}, \code{BACKGROUND}, \code{DEFAULT}) are eligible for
the E-cluster; the upper two are confined to the P-cluster, modulo
thermal pressure. The kernel performs the placement; user code has no
direct affinity API.

This is why \Lucid's data-loader threads
(\chref{ch:data_pipeline}) explicitly leave the QoS class unset,
which inherits \code{USER\_INITIATED} from the calling thread. Setting
it to \code{UTILITY} would have placed I/O workers on E-cores, which
sounds reasonable but is in fact counter-productive: the disk and
network I/O is bound by the SSD's NVMe controller rather than CPU,
and the cost of repeatedly waking an E-core dominates the savings
from keeping the P-cores idle.

\section{Apple Matrix Coprocessor (AMX) and SME}
\label{sec:apple_silicon:amx}

The Apple Matrix Coprocessor is, as of the time of writing, the most
performance-critical and least-documented component of Apple silicon.
It is an undocumented matrix-multiply accelerator shared by the
P-cluster (each P-cluster has one AMX unit; the E-cluster does not)
and accessible only through reserved instruction encodings that are
not part of the architected ARM ISA. Its existence is acknowledged by
Apple only through patent filings and through the throughput of
\Accelerate's BLAS routines, which exceed what is possible with NEON
alone by a factor of four to ten.

\subsection{Instruction format and execution model}
\label{ssec:apple_silicon:amx_isa}

AMX instructions are encoded as ARM-reserved opcodes that the
processor decodes into messages to the shared coprocessor. The unit
operates as a tile machine: it owns a set of internal register tiles
(\code{X}, \code{Y}, \code{Z}) holding submatrices of $32 \times 32$
to $64 \times 64$ depending on element type, and it executes
outer-product micro-operations that accumulate into the \code{Z}
tile. A typical $\code{mlx\_outer}$-style instruction takes two row
vectors from \code{X} and \code{Y} and accumulates their outer
product into a slice of \code{Z}; chaining these reproduces GEMM with
peak throughput around $2$\,TFLOP/s on M1 P-cluster, scaling roughly
linearly with cluster count thereafter.

The encoding has been documented unofficially in reverse-engineering
work by Dougall Johnson and others, and exposed through low-level
libraries such as \texttt{aarch64-jit-amx}. Apple has never blessed
direct AMX programming; the supported path has always been
\Accelerate.

\subsection{Throughput characteristics}
\label{ssec:apple_silicon:amx_perf}

Empirical AMX throughput, measured through \code{cblas\_sgemm}, sits
in the range of $1.5$ to $2.5$ TFLOP/s per P-cluster on
single-precision matrix multiplication of sizes
$\geq 256 \times 256$. Smaller matrices do not amortise the
cluster-internal dispatch cost and run at NEON throughput or lower.
Mixed precision (BF16 inputs, FP32 accumulate) on M3 and later
roughly doubles the peak, but \Accelerate\ has not, as of the 3.5
series of \Lucid, exposed BF16 \code{cblas\_gemm\_bf16bf16f32} on all
SDKs; \Lucid's CPU backend therefore still uses FP32 GEMM by default,
and the BF16 enablement lives in the AMP roadmap.

\subsection{Relationship to ARM SME in the M4 generation}
\label{ssec:apple_silicon:sme}

The M4 introduces the ARM Scalable Matrix Extension (SME) as the
architected successor to AMX. SME is a fully documented ARM ISA
extension (ARMv9.2-A) with portable assembly, intrinsics, and
debugger support. From \Accelerate's perspective the transition is
seamless: \BNNS\ and \code{cblas\_sgemm} produce the same outputs on
SME and AMX, with the dispatcher selecting at runtime based on
\code{sysctl}-reported features.

For \Lucid, three things change with SME:

\begin{enumerate}
  \item The CPU backend (\Accelerate) becomes future-proof: a future
    macOS version that drops AMX shims will still run, because the
    pathway used by \Accelerate\ is now ARM-architected.
  \item The compile path (\chref{ch:compile_design}) acquires a
    portable target. \Lucid\ does not currently emit SME instructions
    directly, but \code{lucid.compile} could in principle emit them
    in a future generation, removing the dependency on \Accelerate's
    in-process dispatch.
  \item The cost model for the CPU backend changes: SME tiles are
    of \emph{variable} width (the \code{SVL} parameter), and the M4
    exposes a tile size of $512$ bits, which makes some smaller-matrix
    operations more efficient than they were on AMX. \Lucid's
    dispatcher does not currently exploit this, but future Accelerate
    versions are expected to.
\end{enumerate}

\section{Apple Neural Engine (ANE)}
\label{sec:apple_silicon:ane}

The Neural Engine is the most marketed and the least useful of Apple
silicon's accelerators from \Lucid's perspective. It is a fixed
function unit: a 16-core (M1, M2) or 16/32-core (M3, M4) NPU
specialised for low-precision neural-network inference, with a peak
throughput on M4 quoted by Apple at $38$ TOPS at INT8.

The ANE is exposed exclusively through CoreML and ML Compute. It does
not accept arbitrary kernels, does not support arbitrary dtypes
(INT8 and FP16 only; no FP32, no BF16 on most generations, no FP64
ever), and operates only on graphs that have been compiled offline
into the CoreML format. There is no SDK that exposes the ANE for
training; there is no online graph submission API.

This places the ANE outside \Lucid's scope by construction. \Lucid\
targets \emph{training} workloads, which require:
\begin{enumerate}
  \item Gradient computation, hence backward kernels for every op.
  \item Mixed precision with FP32 master weights, hence FP32
    arithmetic.
  \item Dynamic graph construction at the Python level, hence online
    submission.
\end{enumerate}
The ANE satisfies none of these. A future inference-only
\code{lucid.deploy} target could compile a trained model to CoreML
and benefit from the ANE on user devices, but the ANE will never be
part of the training graph and is therefore not part of this
chapter's primary subject. We mention it here only to forestall the
common question of why \Lucid\ does not use it.

\section{The GPU}
\label{sec:apple_silicon:gpu}

The GPU is the workhorse. On any non-trivial neural network, the GPU
performs more than $95\,\%$ of the floating-point arithmetic, even
when the CPU side is doing its job (data loading, gradient
accumulation, optimiser updates). \Lucid's compile path
(\chref{ch:graph_capture_wins}) emits its compiled graphs onto the
GPU via \MPSGraph; the \MLX\ runtime executes its eager kernels on
the GPU; the CPU is a coordinator. Understanding the GPU's
microarchitecture is therefore central to understanding \Lucid's
performance profile.

\subsection{Compute hierarchy}
\label{ssec:apple_silicon:gpu_compute}

Apple's GPU is organised as a number of \emph{GPU cores} (Apple's
marketing term; the architectural unit is sometimes called a
\emph{shader core} or \emph{compute unit}). Each GPU core contains:
\begin{itemize}
  \item Several \emph{shader execution units}, each capable of
    executing 32-lane SIMD (\emph{SIMD group} in Metal parlance) FP32
    operations per cycle.
  \item A texture unit and ray-tracing unit (on M3 and later).
  \item Local memory (called \emph{threadgroup memory} in Metal),
    typically $32$\,KB per core.
  \item Register file for in-flight threads (each thread gets a
    private slice).
\end{itemize}

The base M1 has 7 or 8 GPU cores; M1~Pro has 14 or 16; M1~Max has
24 or 32; M1~Ultra has 48 or 64. The M2 and M3 generations followed
a similar scaling pattern, with each tier roughly doubling the
previous. M4~Max as of the time of writing has up to 40 GPU cores.

Each GPU core can sustain in the neighbourhood of $200$ FP32 GFLOP/s
at base clocks. The aggregate throughput thus scales linearly with
core count, modulo memory bandwidth and thermal constraints. The
M1~Max's 32-core GPU delivers approximately $10.4$ TFLOP/s FP32
peak; the M3~Max's 40-core GPU is approximately $14$ TFLOP/s. These
are the numbers that bound \Lucid's training throughput on
matrix-multiply-bound workloads.

\subsection{Memory hierarchy on the GPU side}
\label{ssec:apple_silicon:gpu_mem}

Each GPU core sees memory through a layered hierarchy:
\begin{enumerate}
  \item \textbf{Registers.} Per-thread, allocated by the Metal
    compiler from a fixed pool ($32$\,KB to $64$\,KB per core
    depending on generation). Register pressure governs occupancy.
  \item \textbf{Threadgroup memory.} A small, fast, software-managed
    scratchpad shared by all threads in a threadgroup; typically
    $32$\,KB.
  \item \textbf{Tile memory.} A dedicated on-chip buffer used during
    render-target operations; less relevant for compute.
  \item \textbf{System Level Cache.} Shared with the CPU; the GPU's
    only path to DRAM goes through the SLC.
  \item \textbf{DRAM (LPDDR).} The unified pool.
\end{enumerate}

There is no separate VRAM, no separate L2 the GPU owns exclusively.
The SLC is the only large on-chip GPU cache, and it is shared with
every other client of the fabric. Whether this is an advantage or a
disadvantage depends on the workload: well-behaved single-tenant ML
training benefits from larger effective cache; a system under display
or video-encode pressure can starve the GPU.

\subsection{Dynamic Caching (M3 and later)}
\label{ssec:apple_silicon:gpu_dyncache}

The M3 generation introduced \emph{Dynamic Caching}, which Apple
describes as ``allocating GPU memory in real time based on what each
program needs.'' Concretely, the GPU's register and threadgroup-memory
allocation, which on M1/M2 was determined statically at shader compile
time, becomes runtime-virtualised on M3 and later: the hardware
allocates registers dynamically to in-flight threads, enabling higher
effective occupancy on shaders that previously over-allocated
registers.

For \Lucid\ this is observable but not directly programmable.
Compiled \MLX\ and \MPSGraph\ kernels benefit transparently when run
on M3 or M4 hardware; the same kernels on M1/M2 do not. We document
the empirical effect in \chref{ch:mpsgraph_dispatch}, where the M3
generation shows a step change in per-op throughput for irregular
kernels (reductions over non-contiguous axes, scatter-with-gather
patterns).

\subsection{Hardware ray tracing and mesh shaders}
\label{ssec:apple_silicon:gpu_rt}

The M3 generation also added hardware ray tracing and mesh shader
support. Neither is relevant for ML training and we mention them only
to clarify the silicon budget: a portion of the GPU die area on M3
and later is dedicated to graphics-only features, which is part of
why per-core compute throughput grew less than the marketing core
count would suggest on the M3 transition.

\subsection{Metal Performance Shaders and MPSGraph}
\label{ssec:apple_silicon:gpu_mpsg}

The software interface to the GPU has two layers:
\begin{enumerate}
  \item \textbf{\Metal.} The low-level API: command buffers, shader
    programs in MSL (Metal Shading Language), explicit synchronisation.
    \MLX\ uses Metal directly for its kernels.
  \item \textbf{\MPSGraph.} The high-level API: a graph IR that the
    Metal Performance Shaders framework compiles to optimised
    multi-kernel sequences. \Lucid's compile path
    (\chref{ch:mps_builder}) emits onto \MPSGraph.
\end{enumerate}

These are not mutually exclusive. \MLX\ continues to dispatch most
eager kernels through its own MSL; \MPSGraph\ is used in \Lucid\
only by \code{lucid.compile} and by a small set of per-op
optimisations introduced in 3.4 (\chref{ch:mpsgraph_dispatch}). The
two paths can coexist within a single program because they ultimately
share the same Metal command queue and the same unified memory pool.

\section{Memory System}
\label{sec:apple_silicon:memory}

Every load and store on Apple silicon, regardless of which compute
unit issued it, ultimately resolves through the same memory
controllers to the same LPDDR pool. This section describes the
characteristics of that pool, which set the upper bound on data-bound
workloads (i.e.\ most of deep learning).

\subsection{LPDDR generations and channels}
\label{ssec:apple_silicon:memory_lpddr}

\begin{table}[t]
\centering
\small
\begin{tabular}{lcccc}
\toprule
SoC & DRAM type & Channels & Capacity range & Peak BW (GB/s) \\
\midrule
M1        & LPDDR4X-4266 & 4  & 8--16\,GB     & $\approx 68$    \\
M1 Pro    & LPDDR5-6400  & 4  & 16--32\,GB    & $\approx 200$   \\
M1 Max    & LPDDR5-6400  & 8  & 32--64\,GB    & $\approx 400$   \\
M1 Ultra  & LPDDR5-6400  & 16 & 64--128\,GB   & $\approx 800$   \\
\midrule
M2        & LPDDR5-6400  & 4  & 8--24\,GB     & $\approx 100$   \\
M2 Pro    & LPDDR5-6400  & 4  & 16--32\,GB    & $\approx 200$   \\
M2 Max    & LPDDR5-6400  & 8  & 32--96\,GB    & $\approx 400$   \\
M2 Ultra  & LPDDR5-6400  & 16 & 64--192\,GB   & $\approx 800$   \\
\midrule
M3        & LPDDR5-6400  & 4  & 8--24\,GB     & $\approx 100$   \\
M3 Pro    & LPDDR5-6400  & 3  & 18--36\,GB    & $\approx 150$   \\
M3 Max    & LPDDR5-6400  & 8 or 6 & 36--128\,GB & $\approx 300$--$400$ \\
\midrule
M4        & LPDDR5X-7500 & 4  & 16--32\,GB    & $\approx 120$   \\
M4 Pro    & LPDDR5X-8533 & 4  & 24--48\,GB    & $\approx 273$   \\
M4 Max    & LPDDR5X-8533 & 8  & 36--128\,GB   & $\approx 546$   \\
\bottomrule
\end{tabular}
\caption[Apple silicon memory subsystem summary.]{Memory subsystem
summary across Apple silicon generations. ``Peak BW'' is the
theoretical maximum aggregate bandwidth; sustained bandwidth from a
single client is typically $60$--$80\,\%$ of peak.}
\label{tab:apple_silicon_memory}
\end{table}

\tabref{tab:apple_silicon_memory} summarises the memory subsystem
across generations. Two observations relevant to \Lucid:

\begin{enumerate}
  \item The M3~Pro is the only generation that reduced peak bandwidth
    relative to its predecessor (M2~Pro). This was driven by a
    rebalancing of the die for cost and yield; the practical
    consequence is that some training workloads benchmarked on
    M2~Pro run \emph{slower} on M3~Pro, and \Lucid's performance
    documentation flags this explicitly
    (\chref{ch:stabilization}).
  \item M4~Max's $546$\,GB/s sits between M2~Max and M2~Ultra. For
    \Lucid's GPT-2 base benchmark, this places M4~Max above the
    sustained training throughput of a single M2~Ultra under thermal
    pressure, which is the only documented data-point at the Ultra
    tier; we revisit in \chref{ch:graph_capture_wins}.
\end{enumerate}

\subsection{System Level Cache}
\label{ssec:apple_silicon:memory_slc}

The System Level Cache (SLC) is a last-level cache shared by every
client of the fabric. Its size grows with the chip tier:
$8$\,MB on M1, $24$\,MB on M1~Pro, $48$\,MB on M1~Max, $96$\,MB on
M1~Ultra. M3 and M4 use comparable sizes. SLC is filled and evicted
by hardware-managed coherence; software has no direct allocation
hints other than memory-allocator placement strategies that
\Accelerate\ uses internally.

For ML workloads the SLC affects two things:
\begin{enumerate}
  \item \textbf{Activation reuse.} A forward pass that fits its
    activations into the SLC can be re-read by the backward pass
    without paying DRAM bandwidth twice. The cut-off is generally
    around $50$--$100$ MB of total activations for the Max tier,
    which corresponds to a moderate batch size on ResNet-50 or a
    small batch on a Transformer of GPT-2~base size.
  \item \textbf{Inter-cluster traffic.} When the CPU and GPU exchange
    a buffer through \Lucid's bridge, the SLC absorbs the transfer
    if the buffer is small. This is why our microbenchmarks
    (\chref{ch:perf_methodology}) show roughly the same CPU-to-GPU
    latency for buffers under $4$\,MB regardless of whether the
    buffer was just written or has been resident for some time.
\end{enumerate}

\subsection{Bandwidth versus arithmetic intensity}
\label{ssec:apple_silicon:memory_intensity}

The roofline model captures the trade-off between memory bandwidth
and arithmetic intensity. For an Apple silicon M3~Max with $400$
GB/s of memory bandwidth and roughly $14$\,TFLOP/s of FP32 GPU peak,
the ridge point of the roofline lies at an arithmetic intensity of
$14 \times 10^{12} \,/\, 400 \times 10^9 \approx 35$ FLOPs per byte.
For comparison, an NVIDIA H100 has a ridge point near $250$ FLOPs
per byte: ML kernels that are bandwidth-bound on H100 are often
compute-bound on Apple silicon, and vice versa. This shifts the set
of optimisations that pay off. \Lucid's no-compile sweep
(\chref{ch:no_compile_sweep}) is an example: algorithmic changes
that reduce arithmetic count (using $\mathrm{sum} / N$ in place of
$\mathrm{mean}$ to fold a scalar division earlier) yield real
speedups on Apple silicon precisely because the workload is closer
to compute-bound than memory-bound at typical operation sizes.

\section{The Apple Silicon Family: M1 through M4}
\label{sec:apple_silicon:generations}

Having described the constituent blocks, we now summarise the
generational evolution. Each generation can be understood as a set of
deltas relative to its predecessor: process node, P-core
microarchitecture, GPU enhancements, AMX/SME transition, and memory
subsystem refresh.

\subsection{M1 (2020)}
\label{ssec:apple_silicon:m1}

The M1 introduced the architecture template that all subsequent
generations refined. Salient properties:
\begin{itemize}
  \item TSMC N5; 16 billion transistors.
  \item Firestorm P-cores (8-wide decode, ROB $\approx 630$); 4 of
    them in the base configuration.
  \item Icestorm E-cores (4-wide decode); 4 of them.
  \item GPU: 7 or 8 cores; up to $2.6$\,TFLOP/s FP32.
  \item AMX1.
  \item LPDDR4X-4266, 4 channels, 68\,GB/s, up to 16\,GB.
\end{itemize}
The M1 Pro, Max, and Ultra variants applied the cluster-replication
pattern described earlier: the same blocks, more of them, faster
memory bus.

\subsection{M2 (2022)}
\label{ssec:apple_silicon:m2}

The M2 was an iteration on N5 (specifically, N5P), and applied two
classes of changes:
\begin{itemize}
  \item Avalanche P-cores, Blizzard E-cores. Same width, refined
    schedulers and improved branch prediction; roughly $10$--$15\,\%$
    IPC improvement.
  \item GPU bumped to 10 cores at the base tier; LPDDR transition to
    LPDDR5-6400.
  \item AMX2 (incremental).
\end{itemize}
No fundamental architectural change; this was a refinement
generation.

\subsection{M3 (2023)}
\label{ssec:apple_silicon:m3}

The M3 was a major architectural step:
\begin{itemize}
  \item Move to N3B (3\,nm), a substantial process shrink.
  \item Everest P-cores; further IPC growth.
  \item GPU receives Dynamic Caching, hardware ray tracing, mesh
    shader support. Performance per core grows.
  \item AMX3.
  \item Mixed memory bus configurations: M3~Pro reduced its memory
    bus relative to M2~Pro, an anomaly noted in
    \tabref{tab:apple_silicon_memory}.
\end{itemize}

\subsection{M4 (2024)}
\label{ssec:apple_silicon:m4}

The M4 transitioned to N3E (a refinement of N3) and architecturally
crossed into ARMv9.2-A:
\begin{itemize}
  \item First Apple SoC to expose SME at the ISA level (in place of
    AMX).
  \item LPDDR5X transition (7500 MT/s base; 8533 MT/s on Pro and
    Max).
  \item Continued GPU refinements; mature Dynamic Caching.
  \item Improved Neural Engine quoted at $38$ TOPS.
\end{itemize}
For \Lucid, the M4 transition is the most important since M1: it is
the first generation on which \code{lucid.compile} could plausibly
target SME directly in a future release, removing the dependency on
\Accelerate's in-process dispatch.

\section{Implications for ML Workloads}
\label{sec:apple_silicon:ml_implications}

We close this chapter with the qualitative conclusions that frame the
rest of the book.

\subsection{The dispatch trade-off}
\label{ssec:apple_silicon:ml_dispatch}

Apple silicon's heterogeneous compute means that a framework has at
least three reasonable backends for any given operation: the GPU
(via \MLX\ and/or \MPSGraph), the CPU's NEON pipes (via direct
SIMD), and the CPU's matrix coprocessor (via \Accelerate's \BNNS\
and BLAS). The Neural Engine is a fourth in principle but absent
from training graphs, as discussed.

The performance ordering between these backends is not monotone in
problem size. For small matrices ($\lesssim 64\times 64$), NEON or
even scalar code can be competitive because GPU kernel launch
overhead dominates and AMX's cluster-internal dispatch overhead is
also non-trivial. For medium matrices ($64$--$256$), AMX dominates.
For large matrices ($\geq 512$), the GPU dominates because of its
aggregate throughput advantage. \Lucid's dispatch logic
(\chref{ch:backend_dispatch}) does not attempt to choose between
these dynamically at the framework level: the user picks a device
(\code{cpu} or \code{metal}) and within that device the BLAS or
\MLX\ library applies its own internal heuristic.

\subsection{The memory-bandwidth ceiling}
\label{ssec:apple_silicon:ml_bw}

Apple silicon's most distinctive performance property is its very
high memory bandwidth relative to its compute. On the largest
configurations (Ultra), the bandwidth-to-compute ratio approaches one
byte per FLOP, which is roughly an order of magnitude better than a
discrete GPU. This is what makes large-batch streaming workloads
(autoencoders, certain Transformer training regimes with large
sequence length) practical on Apple silicon despite the modest peak
compute.

The corollary is that bandwidth-bound kernels (layer norm,
softmax, attention, anything that touches every element of a large
tensor exactly once) tend to be the bottleneck on small Apple
silicon configurations. \Lucid's normalisation family rewrites
(\chref{ch:no_compile_sweep}) target these explicitly.

\subsection{The power envelope}
\label{ssec:apple_silicon:ml_power}

Every Apple SoC is bound by a thermal design power (TDP) envelope
that determines sustained performance. The M-series TDPs range from
roughly $20$\,W on the base M1 to roughly $90$\,W on the M1~Ultra;
the M4 family has comparable envelopes. Under sustained load, all
configurations throttle from their peak frequencies, and the GPU
cores are the largest single consumer.

For ML training, this matters in two ways:
\begin{enumerate}
  \item A burst-mode benchmark (the first few iterations of a
    training run) does not predict sustained throughput. \Lucid's
    benchmark methodology (\chref{ch:perf_methodology}) discards
    the first $20$--$50$ iterations to allow the chip to settle.
  \item Active cooling matters. The MacBook Pro chassis throttles
    less aggressively than the MacBook Air; the Mac Studio less than
    either. Numbers in \Lucid's performance documentation are taken
    on Mac Studio M-series hardware to control for this.
\end{enumerate}

\subsection{The compile-vs-eager trade-off}
\label{ssec:apple_silicon:ml_compile}

Finally, the existence of \MPSGraph\ as a graph-level abstraction
above \Metal\ means that an ML framework on Apple silicon faces a
genuine choice between eager dispatch (each op submitted as a Metal
command buffer) and graph capture (a sequence of ops compiled into a
single \MPSGraph\ executable). \MLX\ defaults to eager; \Lucid's
\code{lucid.compile} (\chref{ch:compile_design}) opts into graph
capture for a measurable win on most training workloads. The
trade-offs and the exact wins are documented in
\chref{ch:graph_capture_wins}.

\section{Summary and Forward References}
\label{sec:apple_silicon:summary}

Apple silicon is a tightly integrated, heterogeneous SoC family whose
defining properties --- unified coherent memory, multiple
fixed-function accelerators, and a wide-but-low-frequency CPU complex
--- collectively dictate every interesting design decision in
\Lucid. The five compute units described in this chapter --- P-cores,
E-cores, AMX/SME, the Neural Engine, and the GPU --- form the set of
candidate dispatch targets, of which \Lucid\ uses four (the Neural
Engine is excluded by construction).

The remainder of \partref{part:hardware} expands on the runtime
implications. \chref{ch:unified_memory} formalises the unified memory
model and describes how \Lucid\ exploits it without ever issuing an
explicit copy. \chref{ch:mlx_runtime} introduces \MLX\ as the
GPU-side runtime that \Lucid\ depends on, including its
graph-eager hybrid execution model. \chref{ch:accelerate} introduces
\Accelerate\ as the CPU-side counterpart. \chref{ch:metal_mpsgraph}
treats \Metal\ and \MPSGraph\ directly, since both are visible to
\Lucid's compile path. After these four foundation chapters, the
book proceeds to the system architecture proper in
\partref{part:system}, where the hardware view of this chapter becomes
a software dependency graph that \Lucid\ has been engineered to
respect.

\begin{lucidtip}
For readers primarily interested in framework-level concerns rather
than silicon-level concerns, \secref{sec:apple_silicon:ml_implications}
contains the only consequences they need to internalise: the dispatch
trade-off between three backends, the memory-bandwidth ceiling, the
thermal envelope, and the compile-vs-eager trade-off. The remainder of
this chapter can be returned to as a reference when those concerns
become concrete.
\end{lucidtip}
